\section{Coding: el entorno más general del mundo cibernético}

La sección anterior trató el límite superior de los modelos; esta sección pasa al entorno de acción. Guangmi (广密) llama a Coding «el entorno del mundo cibernético más general» y «las manos del modelo». La frase no afirma que la AGI equivalga a escribir código, sino que el entorno de código tiene tres ventajas: se puede operar, se puede verificar y se puede componer. Cuando el modelo escribe código, lo ejecuta, lee la retroalimentación de las pruebas, llama herramientas y modifica archivos, en esencia está actuando dentro de un mundo digital controlable.

\begin{figure}[H]
\centering
\includegraphics[width=0.96\textwidth]{coding-as-cyber-environment.png}
\caption{Coding es un entorno del mundo cibernético: el código no es una habilidad aislada, sino la puerta de entrada para que el modelo actúe, verifique y transforme el mundo digital. Diagrama conceptual de elaboración propia, a partir de la conversación de 00:11:56--00:19:55.}
\end{figure}

\begin{knowledgebox}{Lectura de la figura: el valor de Coding está en el entorno, no en la sintaxis}
De Code a Tools, Tests, Agent y AGI Path, la figura destaca el ciclo cerrado: el modelo puede producir acciones, las acciones pueden cambiar el entorno y el entorno puede devolver retroalimentación. Frente a las preguntas y respuestas de un chat, en el entorno de código es más fácil definir el éxito y el fracaso; por eso resulta más adecuado para entrenar y evaluar un Agent.
\end{knowledgebox}

\subsection{Por qué Coding es más fundamental que la búsqueda y la recomendación}

Esta subsección explica la analogía de Guangmi. El motor de búsqueda organiza las páginas web como un entorno consultable; el motor de recomendación organiza el contenido como un entorno consumible; Coding, en cambio, organiza la actividad económica digital como un entorno modificable. El modelo no solo encuentra información: también puede crear archivos, llamar API, ejecutar programas, escribir scripts de automatización, probar resultados y corregir errores. Esta propiedad de «poder cambiar el mundo» hace que Coding esté más cerca del cuerpo de un Agent que la simple distribución de información.

\begin{knowledgebox}{Nota de clase: el código es el espacio de acciones del modelo en el mundo digital}
En el lenguaje del aprendizaje por refuerzo, el entorno proporciona estados, las acciones cambian los estados y la retroalimentación evalúa las acciones. El entorno de Coding vuelve muy claros estos elementos: el repo, los archivos, las pruebas, los mensajes de error, la CI, las API, el navegador y la base de datos pueden funcionar como estado y como retroalimentación.
\end{knowledgebox}

\noindent\begin{tabularx}{\textwidth}{p{0.18\textwidth}p{0.34\textwidth}X}
\toprule
Entorno & Acciones principales & Valor para un Agent \\
\midrule
Búsqueda & Consultar, filtrar, resumir & Proporciona información, pero su capacidad de cambiar el mundo es débil. \\
Recomendación & Emparejar, ordenar, consumir & Se orienta más a distribuir tráfico; no es un entorno de acción general. \\
Coding & Leer y escribir código, ejecutar pruebas, llamar herramientas, corregir errores & Se puede operar, verificar y componer; es un entorno de acción digital. \\
Science/Gaming & Experimentar, simular, buscar estrategias & También es sólido, pero por ahora su generalidad y su escala de datos son menores que las de Coding. \\
\bottomrule
\end{tabularx}

\subsection{Computer use y las manos del modelo}

Esta subsección conecta Coding con computer use. Guangmi menciona que Manus hizo que el público percibiera con fuerza, por primera vez, el magic moment de tool use/computer use, y que la capacidad de los modelos de Anthropic es una base importante detrás de ello. Aquí conviene notar una relación de niveles: una empresa de producto puede envolver una capacidad en una experiencia perceptible, pero el modelo subyacente debe tener primero la capacidad de operar herramientas, entender la UI, leer y escribir archivos y mantener el estado en tareas largas.

\begin{importantbox}{La línea divisoria entre «saber decir» y «saber hacer»}
En la era de ChatGPT, los modelos hacían sentir sobre todo que «entienden»; en la era de los Agent, los modelos deben hacer sentir que «pueden completar la tarea». Este cambio requiere herramientas, permisos, memoria, recuperación ante errores y verificación de resultados, no solo respuestas más fluidas.
\end{importantbox}

\subsection{¿Coding se convertirá en la forma final del producto?}

Esta subsección responde a la pregunta de seguimiento de Zhang Xiaojun (张小珺). El juicio de Guangmi se acerca más a esto: Coding es un motor tecnológico, no necesariamente la expresión final del producto. El producto final de un motor de recomendación puede ser un flujo de videos cortos; el producto final de un motor de búsqueda es la caja de búsqueda y la página de respuestas; Coding, como motor, podría manifestarse en el futuro en IDE, automatización de oficina, Agent personales, flujos de trabajo empresariales, cómputo científico y fábricas de software. Para el usuario común, la interfaz final quizá diluya la «programación», pero por dentro seguirá habiendo ejecución de código y de herramientas.

\begin{warningbox}{No hay que equiparar el mercado de herramientas para desarrolladores con el mercado del motor de Coding}
Cursor, Claude Code y los plugins de IDE fueron las primeras formas en despegar, porque los desarrolladores tienen una motivación alta y una retroalimentación clara. Pero si Coding es de verdad el motor general de la economía digital, su producto final no servirá solo a programadores, sino que llegará a la oficina, la investigación, la operación, el diseño, el análisis de negocio y la automatización personal.
\end{warningbox}

\subsection{Resumen de la sección}

La importancia estratégica de Coding está en que da al modelo un mundo digital en el que se puede actuar, verificar y escalar. Es el campo de batalla principal donde los Agent se implementan primero, y también el entorno decisivo para comprobar si un modelo pasa de conversar a actuar.
